% !TEX root = ../lab1_report.tex
% ==========================================================
% 实验二：发动机虚拟结构实验
% ==========================================================
\section{实验二：发动机虚拟结构实验}

\subsection{实验目的}

\begin{enumerate}
    \item 通过实验掌握发动机总体结构。
    \item 掌握压气机、燃烧室、涡轮的组成和结构特点。
    \item 掌握航空发动机附件系统的组成及安装。
\end{enumerate}

\subsection{实验装置}

该实验装置是一台发动机三维电子样机，是一台完全可拆卸的电子样机。实验学生可自主拆卸发动机中任何一个零部件，甚至一个单独转子叶片。

电子样机所展示的发动机为\textbf{某型军用涡扇发动机}，带加力燃烧室，采用双转子构型。

\subsection{实验步骤及内容}

\subsubsection{启动程序}

启动电子样机程序。

\subsubsection{总体结构观察}

观察发动机主体结构，对发动机外形轮廓建立感性认识。

\subsubsection{部件拆卸与观察}

要求拆卸过程中，每一个可拆卸的零件均要拆卸，在拆卸过程中结合结构课内容，观察拆卸部件的结构，根据各部件和各系统的结构和安装形式，分析其作用和工作原理。依次拆卸并观察：

\begin{enumerate}
    \item \textbf{压气机}：低压压气机5级，高压压气机11级（一说12级，确切级数有待确认），共16或17级。观察转子叶片和静子叶片的排布、榫头连接形式、级间密封结构。低压转子转速较低但叶片较长，高压转子叶片短而密集。
    \item \textbf{涡轮}：高压涡轮2级，低压涡轮2级。观察涡轮导向器和转子的叶片形状、冷却结构、榫头形式。
    \item \textbf{燃烧室}：环管燃烧室，共10个火焰筒，环绕发动机轴线周向排列。观察火焰筒内外壁、旋流器、燃油喷嘴、掺混孔和冷却孔的布置。环管燃烧室介于单管燃烧室与环形燃烧室之间，每个火焰筒独立工作，外套共同的环形机匣。
    \item \textbf{加力燃烧室}：观察火焰稳定器、喷油杆和隔热屏的结构。
    \item \textbf{外涵道排气混合器}：观察内外涵气流的掺混方式和混合器构型。
    \item \textbf{防冰系统}：观察从压气机引出热气的管路和防冰腔的结构。
    \item \textbf{传动附件}：观察附件机匣及内部传动齿轮系的布置。
    \item \textbf{主燃油系统}：观察燃油泵、燃油调节器、燃油分配器和喷嘴。
    \item \textbf{滑油系统}：观察滑油箱、滑油泵、油气分离器和回油管路。
    \item \textbf{发动机操纵系统}：观察油门操纵联动机构和作动器。
    \item \textbf{电器设备}：观察点火系统、发电机、电缆布置。
    \item \textbf{防喘调节系统}：观察放气活门和可调静叶机构。
    \item \textbf{加力控制系统}：观察加力燃油泵和喷口面积调节机构。
    \item \textbf{附面层控制系统}：观察引气和吹气装置。
\end{enumerate}

\subsection{实验分析与讨论}

\begin{enumerate}
    \item \textbf{实验中观察到的发动机是一款军用还是民用发动机？它在结构上有哪些特点？}

    观察到的发动机为带加力燃烧室的某型军用涡扇发动机。结构特点包括：采用双转子构型、配备加力燃烧室实现推力增推、环管燃烧室（10个火焰筒）、双涵道设计。

    \item \textbf{压气机的榫头是什么形式？}

    根据实际观察，压气机转子叶片采用燕尾形榫头。燕尾形榫头结构简单，加工方便，通过离心力自锁在轮盘榫槽中，是航空发动机压气机叶片常用的连接方式。

    \item \textbf{燃烧室是什么类型？该燃烧室的基本结构有哪些？}

    本发动机采用环管燃烧室，由10个独立的火焰筒环绕轴线周向排列，外套共同的环形机匣。每个火焰筒的基本结构包括：旋流器（头部进气、产生回流区稳定火焰）、燃油喷嘴（将燃油雾化喷入）、火焰筒壁面（开有掺混孔和冷却孔/气膜冷却环带）。环管燃烧室结合了单管燃烧室（方便单个维修更换）和环形燃烧室（空间利用率高）的优点。10个火焰筒之间通过联焰管连接，保证各火焰筒之间的火焰传播。

    \item \textbf{防冰热气是从第几级压气机引出？}

    防冰热气从高压压气机中间级引出，引气温度和压力需满足防冰要求。引气通过管路送至进气道前缘、进口导叶等易结冰部位进行加热防冰。

    \item \textbf{该发动机的压气机共有几级？涡轮有几级？}

    低压压气机5级，高压压气机11级（一说12级，三维电子样机的级数显示与公开资料存在差异，确切级数有待进一步确认）。高压涡轮2级，低压涡轮2级。该发动机属于中等压比、中等推力的军用涡扇发动机。

    \item \textbf{该发动机的转子支撑方案？}

    根据观察和分析，该发动机的转子支撑方案如下：

    \textbf{高压转子（3个支点）：} 采用0-2-1方案。前支点位于中间机匣，为棍棒轴承（滚柱轴承），承受径向载荷；中间支点位于燃烧室机匣前端，为滚珠轴承，承受轴向推力载荷；后支点位于燃烧室机匣后端，为棍棒轴承，承受径向载荷。

    \textbf{低压转子（4个支点）：} 采用1-2-1方案。前支点位于进气机匣，为滚珠轴承；中间两个支点位于中间机匣区域；后支点位于涡轮后机匣，为棍棒轴承。

    \textbf{中间机匣}：低压转子与高压转子在中间机匣处通过中介轴承（浮环轴承）实现径向定位，低压转子的两个中间支点与高压转子的前支点共用中间机匣结构。

    低压转子4个支点与高压转子3个支点，加上中间机匣处的共用结构，共计约8个承力点。需注意，三维电子样机中部分轴承的具体类型（滚珠/棍棒/浮环）与实际装配可能存在差异，以上分析仅供参考。

    \item \textbf{存在矛盾的若干问题}：
    \begin{enumerate}
        \item 高压压气机级数：电子样机显示为11级，但部分资料记载为12级。两种说法相差1级，可能的解释是电子样机中第一级或末级被合并显示，或不同改型的级数有所调整。
        \item 轴承方案：低压转子4个支点的1-2-1方案与高压转子3个支点的0-2-1方案加中间机匣共8个承力点。但低压转子前支点为滚珠轴承（既承受径向又承受轴向力），与通常低压转子前支点使用棍棒轴承（仅承受径向力）的惯例不一致，需要进一步核实。
        \item 环管燃烧室的10个火焰筒数量与发动机直径和空气流量的匹配关系是否合理，有待结合燃烧室设计手册进一步分析。
    \end{enumerate}
\end{enumerate}

\subsection{结论}

通过本次虚拟结构实验，对该型军用涡扇发动机的总体结构、主要部件（5级低压压气机+11级高压压气机、环管燃烧室、2+2级涡轮）和转子支承方案（低压1-2-1/高压0-2-1）有了较为深入的认识。三维电子样机的可拆卸特性使得学生能够从任意角度观察发动机内部结构，加深了对《航空发动机结构》课程内容的理解。
